\section{Introduction}
\label{sec:introduction}

\subsection{The Redistricting Problem}

Every ten years, following the decennial census, all 50 American states must
redraw their congressional district boundaries.
The process allocates 435 seats across states via the Huntington-Hill method
(Article~I, \S2, operationalised since 1941), then asks each state to partition
its population into exactly $k$ districts where $k$ is its apportioned seat count.
The constitutional floor is simple: districts must be as close to equal in
population as practicable (\wesberry{}, 1964, one person one vote).
Everything above that floor has accumulated as a dense forest of legal requirements
over six decades of redistricting litigation.

The forest has grown denser since 2020.
\rucho{} v.\ \textit{Common Cause} (2019) removed federal courts from partisan
gerrymandering review, but \emph{immediately redirected challengers to state
courts} under state constitutional grounds.
Pennsylvania's \emph{League of Women Voters} (2018) and North Carolina's
\emph{Harper v.\ Hall} (2022) established that state constitutions can ban
partisan gerrymandering even where Rucho closes the federal door.
\allenmilligan{} (2023) reaffirmed and clarified the Voting Rights Act~\S2
obligation to draw minority-opportunity districts where Gingles conditions are met.
And \callais{} (2026) settled the disentanglement question: a state may not
use race-conscious criteria to justify what is actually partisan manipulation,
and the ``strong-inference'' framework requires challengers to show that a
state's stated political goals could have been achieved \emph{with better
minority outcomes}---meaning the algorithm must be capable of separating the
two signals.

This paper provides what the litigation landscape demands:
a \emph{toolbox} of algorithm components where each tool addresses exactly one
legal requirement, and the tools compose without conflict because they are
legally required to be separate.

\subsection{Why a Toolbox, Not a Single Algorithm}

No single redistricting algorithm can satisfy all requirements simultaneously.
The requirements pull in different directions.
Geographic compactness (round districts) conflicts with county preservation
(keeping existing political lines).
VRA minority-opportunity district formation requires concentrating minority
population, which can conflict with equal geographic spread.
Proportionality requires a partisan signal, which \callais{} prohibits from
being mixed with a racial signal.
Census stability across decades requires that the algorithm converges on the
same geographic structure regardless of which census year's data is used.

The toolbox architecture acknowledges these tensions explicitly.
Each tool is a parameter dial:
the \emph{geographic edge weights} of B.2 control compactness;
the \emph{county-sticky weights} of B.10 control subdivision preservation;
the \emph{VRASection alignment score} of B.14 controls Gingles Prong~1 alignment;
the \emph{AreaSection dual constraint} of B.9 controls geographic fairness.
These dials can be tuned independently.
They cannot all be maximised simultaneously, but the tradeoffs are \emph{explicit
and auditable}: a court can read the parameter values from the plan manifest and
understand exactly which tradeoffs were made.

\subsection{The Callais Constraint}

\callais{}, 608 U.S.\ \_\_\_ (2026), has three holdings relevant to algorithmic
redistricting.

First, Louisiana's extra minority district was unconstitutional racial
gerrymandering: the VRA did not require a second majority-minority district
where the first compactness prong of \emph{Thornburg v.\ Gingles} (1986) was
not met.

Second, \S2 challengers must use the ``strong-inference'' framework: they must
show that the state's \emph{stated political goals} (partisan advantage, community
preservation, compactness) could have been achieved \emph{with better minority
outcomes}.
This shifts the burden from ``the map is racially unfair'' to ``a race-neutral
algorithm could have done better on both dimensions''.
The toolbox is precisely positioned to provide this counterfactual: one can
run GeoSection (no minority signal) and VRASection (geographic minority alignment)
on the same state and compare.

Third, at p.~36, the majority holds that race-conscious and partisan signals
must be \emph{disentangled}: a redistricting run that mixes minority VAP
concentrations as a constraint alongside partisan vote share as a constraint
is constitutionally suspect because a court cannot separate the racial motivation
from the partisan motivation in the resulting map.

This third holding is the architecturally decisive constraint.
It means that a redistricting system with multiple simultaneous signals---
minority population fractions plus partisan vote shares in the same METIS
objective function---fails the Callais disentanglement test.
The toolbox's structure, where each mode uses exactly one signal type, is
therefore not a simplification or limitation but a constitutionally correct design.

\subsection{Paper Scope and Contributions}

This paper is the capstone of the B-series algorithmic redistricting programme.
It synthesises findings from B.1 through B.15 into a unified framework.

\paragraph{Contributions.}
\begin{enumerate}
  \item A complete mapping of all eleven legal redistricting requirements to
    specific algorithm components (Section~\ref{sec:requirements-matrix}).
  \item A three-axis taxonomy of algorithm variants (Section~\ref{sec:components}).
  \item A bakeoff of eight algorithm combinations across four litigated states,
    with empirical results on compactness, partisanship, minority representation,
    and county preservation (Section~\ref{sec:bakeoff}).
  \item A formal Callais compliance architecture showing how the toolbox
    enforces p.~36 disentanglement at the CLI level (Section~\ref{sec:callais}).
  \item Practitioner recommendations for which configuration to choose for which
    legal context (Section~\ref{sec:discussion}).
\end{enumerate}

\paragraph{Notation.}
Throughout this paper, $G = (V, E, w)$ denotes the state's census-tract
adjacency graph with vertex population weights and edge boundary-length weights.
$k$ denotes the state's congressional seat count.
$\mathrm{EC}(i:j)$ denotes the minimum edge-cut at split ratio $i:j$.
$\mathrm{EC}_\text{norm}(i:j) = \mathrm{EC}(i:j) / \sqrt{\min(i,j)}$ denotes
the isoperimetrically normalised edge-cut introduced in GeoSection (B.8).
$p^* = L(0.5)$ denotes the Lorenz-curve natural population fraction
(population in the densest 50\,\% of land area).
D\,votes\,\% is the 2020 presidential Democratic two-party vote share.
\end{document}
